% Options for packages loaded elsewhere
\PassOptionsToPackage{unicode}{hyperref}
\PassOptionsToPackage{hyphens}{url}
\PassOptionsToPackage{dvipsnames,svgnames,x11names}{xcolor}
%
\documentclass[
  11pt,
  a4paper]{article}

\usepackage{amsmath,amssymb}
\usepackage{iftex}
\ifPDFTeX
  \usepackage[T1]{fontenc}
  \usepackage[utf8]{inputenc}
  \usepackage{textcomp} % provide euro and other symbols
\else % if luatex or xetex
  \usepackage{unicode-math}
  \defaultfontfeatures{Scale=MatchLowercase}
  \defaultfontfeatures[\rmfamily]{Ligatures=TeX,Scale=1}
\fi
\usepackage{lmodern}
\ifPDFTeX\else  
    % xetex/luatex font selection
\fi
% Use upquote if available, for straight quotes in verbatim environments
\IfFileExists{upquote.sty}{\usepackage{upquote}}{}
\IfFileExists{microtype.sty}{% use microtype if available
  \usepackage[]{microtype}
  \UseMicrotypeSet[protrusion]{basicmath} % disable protrusion for tt fonts
}{}
\makeatletter
\@ifundefined{KOMAClassName}{% if non-KOMA class
  \IfFileExists{parskip.sty}{%
    \usepackage{parskip}
  }{% else
    \setlength{\parindent}{0pt}
    \setlength{\parskip}{6pt plus 2pt minus 1pt}}
}{% if KOMA class
  \KOMAoptions{parskip=half}}
\makeatother
\usepackage{xcolor}
\usepackage[top=30mm,left=20mm,right=20mm,bottom=30mm]{geometry}
\setlength{\emergencystretch}{3em} % prevent overfull lines
\setcounter{secnumdepth}{5}
% Make \paragraph and \subparagraph free-standing
\makeatletter
\ifx\paragraph\undefined\else
  \let\oldparagraph\paragraph
  \renewcommand{\paragraph}{
    \@ifstar
      \xxxParagraphStar
      \xxxParagraphNoStar
  }
  \newcommand{\xxxParagraphStar}[1]{\oldparagraph*{#1}\mbox{}}
  \newcommand{\xxxParagraphNoStar}[1]{\oldparagraph{#1}\mbox{}}
\fi
\ifx\subparagraph\undefined\else
  \let\oldsubparagraph\subparagraph
  \renewcommand{\subparagraph}{
    \@ifstar
      \xxxSubParagraphStar
      \xxxSubParagraphNoStar
  }
  \newcommand{\xxxSubParagraphStar}[1]{\oldsubparagraph*{#1}\mbox{}}
  \newcommand{\xxxSubParagraphNoStar}[1]{\oldsubparagraph{#1}\mbox{}}
\fi
\makeatother

\usepackage{color}
\usepackage{fancyvrb}
\newcommand{\VerbBar}{|}
\newcommand{\VERB}{\Verb[commandchars=\\\{\}]}
\DefineVerbatimEnvironment{Highlighting}{Verbatim}{commandchars=\\\{\}}
% Add ',fontsize=\small' for more characters per line
\usepackage{framed}
\definecolor{shadecolor}{RGB}{241,243,245}
\newenvironment{Shaded}{\begin{snugshade}}{\end{snugshade}}
\newcommand{\AlertTok}[1]{\textcolor[rgb]{0.68,0.00,0.00}{#1}}
\newcommand{\AnnotationTok}[1]{\textcolor[rgb]{0.37,0.37,0.37}{#1}}
\newcommand{\AttributeTok}[1]{\textcolor[rgb]{0.40,0.45,0.13}{#1}}
\newcommand{\BaseNTok}[1]{\textcolor[rgb]{0.68,0.00,0.00}{#1}}
\newcommand{\BuiltInTok}[1]{\textcolor[rgb]{0.00,0.23,0.31}{#1}}
\newcommand{\CharTok}[1]{\textcolor[rgb]{0.13,0.47,0.30}{#1}}
\newcommand{\CommentTok}[1]{\textcolor[rgb]{0.37,0.37,0.37}{#1}}
\newcommand{\CommentVarTok}[1]{\textcolor[rgb]{0.37,0.37,0.37}{\textit{#1}}}
\newcommand{\ConstantTok}[1]{\textcolor[rgb]{0.56,0.35,0.01}{#1}}
\newcommand{\ControlFlowTok}[1]{\textcolor[rgb]{0.00,0.23,0.31}{\textbf{#1}}}
\newcommand{\DataTypeTok}[1]{\textcolor[rgb]{0.68,0.00,0.00}{#1}}
\newcommand{\DecValTok}[1]{\textcolor[rgb]{0.68,0.00,0.00}{#1}}
\newcommand{\DocumentationTok}[1]{\textcolor[rgb]{0.37,0.37,0.37}{\textit{#1}}}
\newcommand{\ErrorTok}[1]{\textcolor[rgb]{0.68,0.00,0.00}{#1}}
\newcommand{\ExtensionTok}[1]{\textcolor[rgb]{0.00,0.23,0.31}{#1}}
\newcommand{\FloatTok}[1]{\textcolor[rgb]{0.68,0.00,0.00}{#1}}
\newcommand{\FunctionTok}[1]{\textcolor[rgb]{0.28,0.35,0.67}{#1}}
\newcommand{\ImportTok}[1]{\textcolor[rgb]{0.00,0.46,0.62}{#1}}
\newcommand{\InformationTok}[1]{\textcolor[rgb]{0.37,0.37,0.37}{#1}}
\newcommand{\KeywordTok}[1]{\textcolor[rgb]{0.00,0.23,0.31}{\textbf{#1}}}
\newcommand{\NormalTok}[1]{\textcolor[rgb]{0.00,0.23,0.31}{#1}}
\newcommand{\OperatorTok}[1]{\textcolor[rgb]{0.37,0.37,0.37}{#1}}
\newcommand{\OtherTok}[1]{\textcolor[rgb]{0.00,0.23,0.31}{#1}}
\newcommand{\PreprocessorTok}[1]{\textcolor[rgb]{0.68,0.00,0.00}{#1}}
\newcommand{\RegionMarkerTok}[1]{\textcolor[rgb]{0.00,0.23,0.31}{#1}}
\newcommand{\SpecialCharTok}[1]{\textcolor[rgb]{0.37,0.37,0.37}{#1}}
\newcommand{\SpecialStringTok}[1]{\textcolor[rgb]{0.13,0.47,0.30}{#1}}
\newcommand{\StringTok}[1]{\textcolor[rgb]{0.13,0.47,0.30}{#1}}
\newcommand{\VariableTok}[1]{\textcolor[rgb]{0.07,0.07,0.07}{#1}}
\newcommand{\VerbatimStringTok}[1]{\textcolor[rgb]{0.13,0.47,0.30}{#1}}
\newcommand{\WarningTok}[1]{\textcolor[rgb]{0.37,0.37,0.37}{\textit{#1}}}

\providecommand{\tightlist}{%
  \setlength{\itemsep}{0pt}\setlength{\parskip}{0pt}}\usepackage{longtable,booktabs,array}
\usepackage{calc} % for calculating minipage widths
% Correct order of tables after \paragraph or \subparagraph
\usepackage{etoolbox}
\makeatletter
\patchcmd\longtable{\par}{\if@noskipsec\mbox{}\fi\par}{}{}
\makeatother
% Allow footnotes in longtable head/foot
\IfFileExists{footnotehyper.sty}{\usepackage{footnotehyper}}{\usepackage{footnote}}
\makesavenoteenv{longtable}
\usepackage{graphicx}
\makeatletter
\def\maxwidth{\ifdim\Gin@nat@width>\linewidth\linewidth\else\Gin@nat@width\fi}
\def\maxheight{\ifdim\Gin@nat@height>\textheight\textheight\else\Gin@nat@height\fi}
\makeatother
% Scale images if necessary, so that they will not overflow the page
% margins by default, and it is still possible to overwrite the defaults
% using explicit options in \includegraphics[width, height, ...]{}
\setkeys{Gin}{width=\maxwidth,height=\maxheight,keepaspectratio}
% Set default figure placement to htbp
\makeatletter
\def\fps@figure{htbp}
\makeatother

\makeatletter
\@ifpackageloaded{caption}{}{\usepackage{caption}}
\AtBeginDocument{%
\ifdefined\contentsname
  \renewcommand*\contentsname{Table of contents}
\else
  \newcommand\contentsname{Table of contents}
\fi
\ifdefined\listfigurename
  \renewcommand*\listfigurename{List of Figures}
\else
  \newcommand\listfigurename{List of Figures}
\fi
\ifdefined\listtablename
  \renewcommand*\listtablename{List of Tables}
\else
  \newcommand\listtablename{List of Tables}
\fi
\ifdefined\figurename
  \renewcommand*\figurename{Figure}
\else
  \newcommand\figurename{Figure}
\fi
\ifdefined\tablename
  \renewcommand*\tablename{Table}
\else
  \newcommand\tablename{Table}
\fi
}
\@ifpackageloaded{float}{}{\usepackage{float}}
\floatstyle{ruled}
\@ifundefined{c@chapter}{\newfloat{codelisting}{h}{lop}}{\newfloat{codelisting}{h}{lop}[chapter]}
\floatname{codelisting}{Listing}
\newcommand*\listoflistings{\listof{codelisting}{List of Listings}}
\makeatother
\makeatletter
\makeatother
\makeatletter
\@ifpackageloaded{caption}{}{\usepackage{caption}}
\@ifpackageloaded{subcaption}{}{\usepackage{subcaption}}
\makeatother

\ifLuaTeX
  \usepackage{selnolig}  % disable illegal ligatures
\fi
\usepackage{bookmark}

\IfFileExists{xurl.sty}{\usepackage{xurl}}{} % add URL line breaks if available
\urlstyle{same} % disable monospaced font for URLs
\hypersetup{
  pdftitle={Data Manipulation with the R package dplyr},
  pdfauthor={Ashish Mandlik},
  colorlinks=true,
  linkcolor={blue},
  filecolor={Maroon},
  citecolor={Blue},
  urlcolor={Blue},
  pdfcreator={LaTeX via pandoc}}


\title{Data Manipulation with the R package dplyr}
\usepackage{etoolbox}
\makeatletter
\providecommand{\subtitle}[1]{% add subtitle to \maketitle
  \apptocmd{\@title}{\par {\large #1 \par}}{}{}
}
\makeatother
\subtitle{Data Analysis for Decision-Making}
\author{Ashish Mandlik}
\date{2024-12-03}

\begin{document}
\maketitle
\begin{abstract}
This presentation explores the use of the dplyr package in R for
efficient data manipulation and analysis. Focused on a practical
scenario, we will demonstrate how to identify the best employees for a
new project by leveraging dplyr's powerful functions. Participants will
learn to filter data based on specific criteria, arrange it for ranking,
and select key information to support decision-making.
\end{abstract}


\section{Introduction}\label{introduction}

R, one of the most popular programming languages for statistical
analysis and data science, offers powerful tools for handling
data---among them, the dplyr package stands out for its simplicity and
versatility.

Data Manipulation with dplyr in R involves transforming and analyzing
data efficiently using a set of functions designed for data wrangling.
dplyr is part of the tidyverse package and provides a coherent grammar
for data manipulation, which helps streamline common data operations.

\section{Types of Verbs}\label{types-of-verbs}

\textbf{Rows}

The most important verbs that operate on rows of a dataset are
\textbf{filter()}, which changes which rows are present without changing
their order, and \textbf{arrange()}, which changes the order of the rows
without changing which are present. Both functions only affect the rows,
and the columns are left unchanged. We'll also discuss
\textbf{distinct()} which finds rows with unique values. Unlike
arrange() and filter() it can also optionally modify the columns.

\textbf{Columns}

There are four important verbs that affect the columns without changing
the rows:\textbf{mutate()}~creates new columns that are derived from the
existing columns,\textbf{select()}~changes which columns are
present,\textbf{rename()}~changes the names of the columns, and
\textbf{relocate()}~changes the positions of the columns.

\phantomsection\label{tbl-tpc}
\begin{longtable}[]{@{}
  >{\raggedright\arraybackslash}p{(\columnwidth - 4\tabcolsep) * \real{0.2500}}
  >{\raggedright\arraybackslash}p{(\columnwidth - 4\tabcolsep) * \real{0.3750}}
  >{\raggedright\arraybackslash}p{(\columnwidth - 4\tabcolsep) * \real{0.3750}}@{}}
\toprule\noalign{}
\begin{minipage}[b]{\linewidth}\raggedright
Sr.no
\end{minipage} & \begin{minipage}[b]{\linewidth}\raggedright
Verbs
\end{minipage} & \begin{minipage}[b]{\linewidth}\raggedright
Function
\end{minipage} \\
\midrule\noalign{}
\endhead
\bottomrule\noalign{}
\endlastfoot
1 & filter() & picks cases based on their values. \\
2 & arrange() & changes the ordering of the rows. \\
3 & distance() & This function calculates a variety of dissimilarity or
distance metrics. \\
4 & mutate() & adds new variables that are functions of existing
variables \\
5 & select() & picks variables based on their names. \\
6 & rename() & changes the names of individual variables \\
7 & relocate() & sed to rearrange the order of columns within a data
frame. \\
8 & summarise() & reduces multiple values down to a single summary. \\
& & \\
\end{longtable}

\textbf{Pipe(} \textbf{\%\textgreater\%)}

Pipe in R~is an operator that takes the output of one function and
passes it into another function as an argument. It links together all
the steps in data analysis making the code more efficient and readable.

Data analysis often involves many steps. A typical journey from raw data
to results might involve filtering cases, transforming values,
summarizing data and then running a~statistical test. Pipe in R links
all these steps together, while keeping our code efficient and readable.

Syntex :

mtcars \textbf{\%\textgreater\%} group\_by(cyl) \%\textgreater\%
summarise(mean\_mpg = mean(mpg))

\section{Exercise}\label{exercise}

\subsection{\texorpdfstring{\textbf{Installation of dplyr
package}}{Installation of dplyr package}}\label{installation-of-dplyr-package}

Open your R studio and in the Console just write this code. R will
automatically install the package.

\begin{verbatim}
install.packages("dplyr")
\end{verbatim}

\textbf{Scenario}

The company needs to select the best employees for a new tech project.
The project requires:

\begin{enumerate}
\def\labelenumi{\arabic{enumi}.}
\item
  \textbf{Relevant roles}: Software Engineers, Data Scientists, DevOps
  Engineers, or Front-End Developers.
\item
  \textbf{Years of experience}: At least 5 years.
\item
  \textbf{Key skills}: Python, AWS, or JavaScript.
\end{enumerate}

The goal is to use the dplyr package to filter and rank employees based
on the given criteria.

\subsection{Steps for the Exercise}\label{steps-for-the-exercise}

Step1: Load and Inspect the Dataset

Start by loading the dataset into R

\begin{Shaded}
\begin{Highlighting}[]
\CommentTok{\# Create the dataset}
\NormalTok{employee\_data }\OtherTok{\textless{}{-}} \FunctionTok{data.frame}\NormalTok{(}
  \AttributeTok{Name =} \FunctionTok{c}\NormalTok{(}\StringTok{"Alice Johnson"}\NormalTok{, }\StringTok{"Bob Smith"}\NormalTok{, }\StringTok{"Charlie Brown"}\NormalTok{, }\StringTok{"Dana White"}\NormalTok{, }\StringTok{"Eve Jackson"}\NormalTok{,}
           \StringTok{"Frank Green"}\NormalTok{, }\StringTok{"Grace Lee"}\NormalTok{, }\StringTok{"Henry Wilson"}\NormalTok{, }\StringTok{"Ian Black"}\NormalTok{, }\StringTok{"Jane Doe"}\NormalTok{),}
  \AttributeTok{Role =} \FunctionTok{c}\NormalTok{(}\StringTok{"Software Engineer"}\NormalTok{, }\StringTok{"Data Scientist"}\NormalTok{, }\StringTok{"Project Manager"}\NormalTok{, }\StringTok{"UX Designer"}\NormalTok{,}
           \StringTok{"DevOps Engineer"}\NormalTok{, }\StringTok{"Business Analyst"}\NormalTok{, }\StringTok{"Front{-}End Developer"}\NormalTok{, }
           \StringTok{"Cybersecurity Analyst"}\NormalTok{, }\StringTok{"Product Manager"}\NormalTok{, }\StringTok{"Marketing Specialist"}\NormalTok{),}
  \AttributeTok{Years\_of\_Experience =} \FunctionTok{c}\NormalTok{(}\DecValTok{5}\NormalTok{, }\DecValTok{7}\NormalTok{, }\DecValTok{10}\NormalTok{, }\DecValTok{4}\NormalTok{, }\DecValTok{6}\NormalTok{, }\DecValTok{8}\NormalTok{, }\DecValTok{3}\NormalTok{, }\DecValTok{9}\NormalTok{, }\DecValTok{12}\NormalTok{, }\DecValTok{5}\NormalTok{),}
  \AttributeTok{Skill\_Sets =} \FunctionTok{c}\NormalTok{(}\StringTok{"Python, Java, SQL"}\NormalTok{, }\StringTok{"R, Python, Machine Learning"}\NormalTok{, }\StringTok{"Agile, Scrum, Leadership"}\NormalTok{, }
                 \StringTok{"Figma, Sketch, User Research"}\NormalTok{, }\StringTok{"AWS, Docker, Kubernetes"}\NormalTok{, }\StringTok{"Excel, SQL, Power BI"}\NormalTok{, }
                 \StringTok{"HTML, CSS, JavaScript, React"}\NormalTok{, }\StringTok{"Network Security, Ethical Hacking"}\NormalTok{, }
                 \StringTok{"Product Strategy, Market Analysis"}\NormalTok{, }\StringTok{"SEO, Content Marketing, Analytics"}\NormalTok{),}
  \AttributeTok{Age =} \FunctionTok{c}\NormalTok{(}\DecValTok{28}\NormalTok{, }\DecValTok{32}\NormalTok{, }\DecValTok{40}\NormalTok{, }\DecValTok{27}\NormalTok{, }\DecValTok{30}\NormalTok{, }\DecValTok{35}\NormalTok{, }\DecValTok{25}\NormalTok{, }\DecValTok{36}\NormalTok{, }\DecValTok{38}\NormalTok{, }\DecValTok{29}\NormalTok{)}
\NormalTok{)}

\CommentTok{\# View the data}
\FunctionTok{head}\NormalTok{(employee\_data)}
\end{Highlighting}
\end{Shaded}

Step 2: Filter Employees Based on Role and Experience

Use filter() to select employees with the relevant roles and at least 5
years of experience.

\begin{Shaded}
\begin{Highlighting}[]
\CommentTok{\# Filter relevant roles and experience}
\NormalTok{filtered\_employees }\OtherTok{\textless{}{-}}\NormalTok{ employee\_data }\SpecialCharTok{\%\textgreater{}\%}
  \FunctionTok{filter}\NormalTok{(Role }\SpecialCharTok{\%in\%} \FunctionTok{c}\NormalTok{(}\StringTok{"Software Engineer"}\NormalTok{, }\StringTok{"Data Scientist"}\NormalTok{, }\StringTok{"DevOps Engineer"}\NormalTok{, }\StringTok{"Front{-}End Developer"}\NormalTok{) }\SpecialCharTok{\&}
\NormalTok{         Years\_of\_Experience }\SpecialCharTok{\textgreater{}=} \DecValTok{5}\NormalTok{)}

\NormalTok{filtered\_employees}
\end{Highlighting}
\end{Shaded}

Step 3: Check for Key Skills

Next,filter employees who have the required skills: Python, AWS, or
JavaScript.

\begin{Shaded}
\begin{Highlighting}[]
\CommentTok{\# Filter employees with key skills}
\NormalTok{best\_candidates }\OtherTok{\textless{}{-}}\NormalTok{ filtered\_employees }\SpecialCharTok{\%\textgreater{}\%}
  \FunctionTok{filter}\NormalTok{(}\FunctionTok{grepl}\NormalTok{(}\StringTok{"Python|AWS|JavaScript"}\NormalTok{, Skill\_Sets))}

\NormalTok{best\_candidates}
\end{Highlighting}
\end{Shaded}

Step 4: Rank Candidates by Experience

Use arrange() to rank the candidates by their years of experience in
descending order.

\begin{Shaded}
\begin{Highlighting}[]
\CommentTok{\# Rank candidates by experience}
\NormalTok{ranked\_candidates }\OtherTok{\textless{}{-}}\NormalTok{ best\_candidates }\SpecialCharTok{\%\textgreater{}\%}
  \FunctionTok{arrange}\NormalTok{(}\FunctionTok{desc}\NormalTok{(Years\_of\_Experience))}

\NormalTok{ranked\_candidates}
\end{Highlighting}
\end{Shaded}

Step 5: Select Key Information

Use select() to display only the most relevant columns for
decision-making.

\begin{Shaded}
\begin{Highlighting}[]
\NormalTok{Copy code}
\CommentTok{\# Select relevant columns}
\NormalTok{final\_selection }\OtherTok{\textless{}{-}}\NormalTok{ ranked\_candidates }\SpecialCharTok{\%\textgreater{}\%}
  \FunctionTok{select}\NormalTok{(Name, Role, Years\_of\_Experience, Skill\_Sets)}

\NormalTok{final\_selection}
\end{Highlighting}
\end{Shaded}

\section{Conclusion}\label{conclusion}

In this session, we've seen how the \texttt{dplyr} package in R
simplifies data manipulation through its intuitive and powerful
functions. From filtering and arranging data to summarizing key
insights, \texttt{dplyr} provides a streamlined workflow that enhances
both efficiency and readability.

Using a practical example, we demonstrated how to identify the best
employees for a project, showcasing the value of tools like
\texttt{filter()}, \texttt{arrange()}, and \texttt{mutate()} in
real-world scenarios. These skills are essential for anyone working with
data, enabling faster decision-making and more insightful analysis.

As you continue to explore \texttt{dplyr}, you'll discover even more
ways to handle complex data tasks with ease.

\section{References}\label{references}

\begin{itemize}
\item
  Wickham, H., François, R., Henry, L., \& Müller, K. (2023).
  \emph{dplyr: A Grammar of Data Manipulation}. Retrieved from
  \href{https://dplyr.tidyverse.org/}{https://dplyr.tidyverse.org}
\item
  Tan, P.-N., Steinbach, M., \& Kumar, V. (2018). \emph{Introduction to
  data mining} (2nd ed.). Pearson. Wickham, H., \& Grolemund, G. (2023).
  \emph{R for data science (2e)}. \url{https://r4ds.hadley.nz/}
\end{itemize}




\end{document}
